
\section{数据分析与结果讨论}


\subsection{具体数据}


\subsubsection{主问题数据}

主问题数据主要来自附件一outputs中 four case framework 和 mode decision 部分，并将其输出结果整合为了附件二中的 01--13 号工作表，并在附件三中整理了单独运行调试结果。其中，01--08 号工作表对应 CP-SAT 求解器、融合剪枝的隐枚举法、启发式算法和热启动后的隐枚举法在二臂和三臂场景下的单独调度结果；09--12 号工作表对应四种方法的二臂/三臂时间能耗对比结果；13 号工作表对应二臂/三臂模式决策汇总。

在实验设置上，本文对 $16$ 种任务数量组合和 $3$ 个随机 seed 进行批量实验。每一种任务数量组合用 \texttt{counts\_code} 表示，例如 \texttt{1111} 表示 Type1、Type2、Type3、Type4 四类任务各有 $1$ 个，\texttt{2222} 表示四类任务各有 $2$ 个。不同 seed 表示同一任务规模下不同的物块随机空间分布，因此能够反映任务位置变化对调度结果的影响。

二臂/三臂模式决策的具体结果如表 \ref{tab:mode_decision_16groups} 所示。表中每一行对应一种任务数量组合，数据为同一 \texttt{counts\_code} 下 $3$ 个随机 seed 的平均结果。
```latex
四种方法在 $16$ 种任务组合下的平均结果分别如表 \ref{tab:cp_sat_16groups}--表 \ref{tab:warmstart_16groups} 所示。每一行对应一种 Type 任务数量组合，表中数值为该组合下 $3$ 个随机 seed 的平均值。

\begin{table}[H]
	\centering
	\caption{CP-SAT求解器在16种任务组合下的平均结果}
	\label{tab:cp_sat_16groups}
	\tiny
	\setlength{\tabcolsep}{2.5pt}
	\resizebox{\textwidth}{!}{
		\begin{tabular}{ccccccccc}
			\hline
			Type1 & Type2 & Type3 & Type4 & $C_{\max}^{\mathrm{2B}}$ & $E_{\mathrm{tot}}^{\mathrm{2B}}$ & $C_{\max}^{\mathrm{3B}}$ & $E_{\mathrm{tot}}^{\mathrm{3B}}$ & 平均计算时间/s \\
			\hline
			1 & 1 & 1 & 1 & 1864.00 & 8015.00 & 1144.33 & 8445.00 & 1.1087 \\
			1 & 1 & 1 & 2 & 2390.33 & 12076.33 & 1674.67 & 12578.00 & 0.4266 \\
			1 & 1 & 2 & 1 & 2341.00 & 9844.00 & 1646.67 & 10282.33 & 0.3741 \\
			1 & 1 & 2 & 2 & 2864.67 & 13975.67 & 2172.67 & 14493.33 & 2.8354 \\
			1 & 2 & 1 & 1 & 2160.33 & 8401.00 & 1387.00 & 8807.67 & 0.3500 \\
			1 & 2 & 1 & 2 & 2689.33 & 12514.67 & 1667.33 & 12997.00 & 0.9333 \\
			1 & 2 & 2 & 1 & 2799.67 & 11139.67 & 1842.00 & 11610.33 & 1.6254 \\
			1 & 2 & 2 & 2 & 3300.67 & 15014.33 & 2285.00 & 15531.33 & 17.3515 \\
			2 & 1 & 1 & 1 & 1955.33 & 7985.67 & 1303.33 & 8454.00 & 0.3118 \\
			2 & 1 & 1 & 2 & 2433.00 & 11829.33 & 1644.33 & 12374.00 & 0.9109 \\
			2 & 1 & 2 & 1 & 2573.00 & 10464.67 & 1746.00 & 10947.67 & 1.0879 \\
			2 & 1 & 2 & 2 & 3082.00 & 14612.00 & 2295.33 & 15207.00 & 19.1113 \\
			2 & 2 & 1 & 1 & 2077.67 & 8824.67 & 1505.33 & 9260.33 & 0.6799 \\
			2 & 2 & 1 & 2 & 2638.33 & 13213.00 & 1794.00 & 13787.33 & 3.2290 \\
			2 & 2 & 2 & 1 & 2572.33 & 10689.67 & 1749.00 & 11194.67 & 3.6314 \\
			2 & 2 & 2 & 2 & 3125.00 & 15068.33 & 2282.00 & 15655.00 & 62.2421 \\
			\hline
		\end{tabular}
	}
\end{table}

\begin{table}[H]
	\centering
	\caption{融合剪枝的隐枚举法在16种任务组合下的平均结果}
	\label{tab:pruned_enum_16groups}
	\tiny
	\setlength{\tabcolsep}{2.5pt}
	\resizebox{\textwidth}{!}{
		\begin{tabular}{ccccccccc}
			\hline
			Type1 & Type2 & Type3 & Type4 & $C_{\max}^{\mathrm{2B}}$ & $E_{\mathrm{tot}}^{\mathrm{2B}}$ & $C_{\max}^{\mathrm{3B}}$ & $E_{\mathrm{tot}}^{\mathrm{3B}}$ & 平均计算时间/s \\
			\hline
			1 & 1 & 1 & 1 & 1864.00 & 8015.00 & 1144.33 & 8445.00 & 0.0160 \\
			1 & 1 & 1 & 2 & 2390.33 & 12076.33 & 1674.67 & 12578.00 & 0.1280 \\
			1 & 1 & 2 & 1 & 2341.00 & 9844.00 & 1646.67 & 10282.33 & 0.1290 \\
			1 & 1 & 2 & 2 & 2864.67 & 13975.67 & 2172.67 & 14493.33 & 1.7181 \\
			1 & 2 & 1 & 1 & 2160.33 & 8401.00 & 1387.00 & 8807.67 & 0.1129 \\
			1 & 2 & 1 & 2 & 2689.33 & 12514.67 & 1667.33 & 12997.00 & 0.5551 \\
			1 & 2 & 2 & 1 & 2799.67 & 11139.67 & 1842.00 & 11610.33 & 0.7532 \\
			1 & 2 & 2 & 2 & 3300.67 & 15014.33 & 2285.00 & 15531.33 & 18.3875 \\
			2 & 1 & 1 & 1 & 1955.33 & 7985.67 & 1303.33 & 8454.00 & 0.1284 \\
			2 & 1 & 1 & 2 & 2433.00 & 11829.33 & 1644.33 & 12374.00 & 0.7779 \\
			2 & 1 & 2 & 1 & 2573.00 & 10464.67 & 1746.00 & 10947.67 & 0.8380 \\
			2 & 1 & 2 & 2 & 3082.00 & 14612.00 & 2295.33 & 15207.00 & 27.9397 \\
			2 & 2 & 1 & 1 & 2077.67 & 8824.67 & 1505.33 & 9260.33 & 1.2688 \\
			2 & 2 & 1 & 2 & 2638.33 & 13213.00 & 1794.00 & 13787.33 & 3.3501 \\
			2 & 2 & 2 & 1 & 2572.33 & 10689.67 & 1749.00 & 11194.67 & 4.2431 \\
			2 & 2 & 2 & 2 & 3125.00 & 15068.33 & 2282.00 & 15655.00 & 94.0316 \\
			\hline
		\end{tabular}
	}
\end{table}

\begin{table}[H]
	\centering
	\caption{启发式算法在16种任务组合下的平均结果}
	\label{tab:heuristic_16groups}
	\tiny
	\setlength{\tabcolsep}{2.5pt}
	\resizebox{\textwidth}{!}{
		\begin{tabular}{ccccccccc}
			\hline
			Type1 & Type2 & Type3 & Type4 & $C_{\max}^{\mathrm{2B}}$ & $E_{\mathrm{tot}}^{\mathrm{2B}}$ & $C_{\max}^{\mathrm{3B}}$ & $E_{\mathrm{tot}}^{\mathrm{3B}}$ & 平均计算时间/s   \\
			\hline
			1 & 1 & 1 & 1 & 1864.00 & 8015.00 & 1144.33 & 8445.00 & 0.1873  \\
			1 & 1 & 1 & 2 & 2390.33 & 12076.33 & 1674.67 & 12578.00 & 0.2025  \\
			1 & 1 & 2 & 1 & 2341.00 & 9844.00 & 1646.67 & 10289.00 & 0.2024  \\
			1 & 1 & 2 & 2 & 2864.67 & 13975.67 & 2172.67 & 14493.33 & 1.0874  \\
			1 & 2 & 1 & 1 & 2160.33 & 8401.00 & 1398.00 & 8833.67 & 0.2134  \\
			1 & 2 & 1 & 2 & 2689.33 & 12514.67 & 1676.33 & 13005.33 & 1.0964  \\
			1 & 2 & 2 & 1 & 2804.00 & 11122.33 & 1842.00 & 11610.33 & 1.1362  \\
			1 & 2 & 2 & 2 & 3305.00 & 14997.00 & 2286.00 & 15533.33 & 1.0779  \\
			2 & 1 & 1 & 1 & 1955.33 & 7985.67 & 1338.33 & 8482.00 & 0.2060  \\
			2 & 1 & 1 & 2 & 2433.00 & 11829.33 & 1654.33 & 12395.33 & 1.0458  \\
			2 & 1 & 2 & 1 & 2573.00 & 10464.67 & 1749.00 & 10945.00 & 1.1137 \\
			2 & 1 & 2 & 2 & 3082.00 & 14612.00 & 2298.00 & 15211.33 & 1.2413 \\
			2 & 2 & 1 & 1 & 2079.33 & 8815.67 & 1535.67 & 9255.67 & 1.1127 \\
			2 & 2 & 1 & 2 & 2638.33 & 13213.00 & 1967.67 & 13797.33 & 1.2333 \\
			2 & 2 & 2 & 1 & 2572.33 & 10689.67 & 1840.00 & 11189.33 & 1.0568  \\
			2 & 2 & 2 & 2 & 3131.67 & 15072.67 & 2282.67 & 15678.67 & 1.8600  \\
			\hline
		\end{tabular}
	}
\end{table}

\begin{table}[H]
	\centering
	\caption{热启动后的隐枚举法在16种任务组合下的平均结果}
	\label{tab:warmstart_16groups}
	\tiny
	\setlength{\tabcolsep}{2.5pt}
	\resizebox{\textwidth}{!}{
		\begin{tabular}{ccccccccc}
			\hline
			Type1 & Type2 & Type3 & Type4 & $C_{\max}^{\mathrm{2B}}$ & $E_{\mathrm{tot}}^{\mathrm{2B}}$ & $C_{\max}^{\mathrm{3B}}$ & $E_{\mathrm{tot}}^{\mathrm{3B}}$ & 平均计算时间/s \\
			\hline
			1 & 1 & 1 & 1 & 1864.00 & 8015.00 & 1144.33 & 8445.00 & 0.0374 \\
			1 & 1 & 1 & 2 & 2390.33 & 12076.33 & 1674.67 & 12578.00 & 0.1635 \\
			1 & 1 & 2 & 1 & 2341.00 & 9844.00 & 1646.67 & 10282.33 & 0.1564 \\
			1 & 1 & 2 & 2 & 2864.67 & 13975.67 & 2172.67 & 14493.33 & 1.6760 \\
			1 & 2 & 1 & 1 & 2160.33 & 8401.00 & 1387.00 & 8807.67 & 0.1565 \\
			1 & 2 & 1 & 2 & 2689.33 & 12514.67 & 1667.33 & 12997.00 & 0.6132 \\
			1 & 2 & 2 & 1 & 2799.67 & 11139.67 & 1842.00 & 11610.33 & 0.7846 \\
			1 & 2 & 2 & 2 & 3300.67 & 15014.33 & 2285.00 & 15531.33 & 9.8419 \\
			2 & 1 & 1 & 1 & 1955.33 & 7985.67 & 1303.33 & 8454.00 & 0.1477 \\
			2 & 1 & 1 & 2 & 2433.00 & 11829.33 & 1644.33 & 12374.00 & 0.7320 \\
			2 & 1 & 2 & 1 & 2573.00 & 10464.67 & 1746.00 & 10947.67 & 0.8028 \\
			2 & 1 & 2 & 2 & 3082.00 & 14612.00 & 2295.33 & 15207.00 & 23.5496 \\
			2 & 2 & 1 & 1 & 2077.67 & 8824.67 & 1505.33 & 9260.33 & 0.5756 \\
			2 & 2 & 1 & 2 & 2638.33 & 13213.00 & 1794.00 & 13787.33 & 4.1532 \\
			2 & 2 & 2 & 1 & 2572.33 & 10689.67 & 1749.00 & 11194.67 & 4.7159 \\
			2 & 2 & 2 & 2 & 3125.00 & 15068.33 & 2282.00 & 15655.00 & 48.9079 \\
			\hline
		\end{tabular}
	}
\end{table}


\subsubsection{延伸问题数据}


主问题数据主要来自附件一outputs中 nonlinear programming 部分，并将其输出结果整合为了附件二中的 14-17 号工作表，问题一为速度--能耗非线性优化问题，数据来自15 调速能耗优化和14 调速KKT验证；问题二为鲁棒速度模式选择问题，数据来自17鲁棒调速模式选择”和16 鲁棒KKT验证。

由于原始 Excel 中问题一输出结果共有 $1024$ 行，问题二输出结果共有 $256$ 行，若全部列入正文会造成篇幅过长。因此，本文采用“整体统计 + 代表性抽样”的方式展示数据：整体统计用于说明总体趋势，代表性抽样用于展示具体任务组合、截止时间比例、速度倍率、能耗和推荐模式之间的关系。

\paragraph{问题一：速度--能耗非线性优化输出结果}

问题一对同一任务数量组合下多个 seed 的主问题结果取平均，并在平均结果基础上优化单臂任务速度倍率和双臂协同任务速度倍率。该问题共考虑 $16$ 种任务数量组合、$16$ 个截止时间比例和 $4$ 组速度--能耗参数，因此主结果共有
\[
16\times 16\times 4=1024
\]
行。其整体输出结果如表 \ref{tab:problem1_output_overall} 所示。

\begin{table}[H]
	\centering
	\caption{问题一速度--能耗非线性优化整体输出统计}
	\label{tab:problem1_output_overall}
	\footnotesize
	\setlength{\tabcolsep}{4pt}
	\begin{tabular}{lcccccc}
		\hline
		数据范围 & 主结果行数 & 推荐二臂 & 推荐三臂 & 三臂比例 & 二臂平均能耗 & 三臂平均能耗 \\
		\hline
		全部参数组 & 1024 & 795 & 229 & 22.36\% & 11675.21 & 11827.55 \\
		基准参数组 & 256 & 209 & 47 & 18.36\% & 11690.09 & 11948.25 \\
		\hline
	\end{tabular}
\end{table}

为了进一步展示具体计算结果，表 \ref{tab:problem1_output_samples} 选取了若干具有代表性的任务组合和截止时间比例。表中 $s_S$ 表示单臂任务速度倍率，$s_D$ 表示双臂协同任务速度倍率，$E^{2B*}$ 和 $E^{3B*}$ 分别表示二臂、三臂调速后的优化能耗。

\begin{table}[H]
	\centering
	\caption{问题一速度--能耗非线性优化代表性输出样例}
	\label{tab:problem1_output_samples}
	\tiny
	\setlength{\tabcolsep}{2.5pt}
	\resizebox{\textwidth}{!}{
		\begin{tabular}{lcccccccccc}
			\hline
			参数组 & 任务组合 & 截止比例 & 截止时间 
			& 二臂速度 $(s_S,s_D)$ & $E^{2B*}$ & $C_{\max}^{2B*}$ 
			& 三臂速度 $(s_S,s_D)$ & $E^{3B*}$ & $C_{\max}^{3B*}$ 
			& 推荐模式 \\
			\hline
			基准参数 & \texttt{1221} & 1.00 & 2799.67 
			& $(0.9795,1.0214)$ & 11136.16 & 2799.67
			& $(0.9163,0.9687)$ & 11579.97 & 1955.90
			& 二臂 \\
			低慢速惩罚 & \texttt{2122} & 0.70 & 2157.40
			& $(1.4413,1.4192)$ & 17517.08 & 2157.40
			& $(1.0556,1.0703)$ & 15486.80 & 2157.40
			& 三臂 \\
			低慢速惩罚 & \texttt{2122} & 0.80 & 2465.60
			& $(1.2539,1.2471)$ & 16012.73 & 2465.60
			& $(0.9107,0.9468)$ & 14982.40 & 2465.60
			& 三臂 \\
			低慢速惩罚 & \texttt{2221} & 0.70 & 1800.63
			& $(1.4380,1.4162)$ & 12813.07 & 1800.63
			& $(0.9594,0.9877)$ & 11115.36 & 1800.63
			& 三臂 \\
			基准参数 & \texttt{2222} & 0.70 & 2187.50
			& $(1.4314,1.4257)$ & 16984.85 & 2187.50
			& $(1.0259,1.0611)$ & 15720.41 & 2187.50
			& 三臂 \\
			\hline
		\end{tabular}
	}
\end{table}

由表 \ref{tab:problem1_output_samples} 可以看出，推荐模式会随任务组合、截止时间比例和速度--能耗参数发生变化。例如，在 \texttt{1221} 且截止比例为 $1.00$ 的基准参数组下，二臂优化后能耗低于三臂，因此推荐二臂；而在 \texttt{2122}、\texttt{2221} 和 \texttt{2222} 的紧工期样例中，三臂虽然需要承担额外系统能耗，但由于并行能力更强、所需速度倍率较低，调速后能耗反而低于二臂，因此推荐三臂。这说明问题一的结果不是简单地固定推荐某一种模式，而是由任务结构、工期约束和速度--能耗曲线共同决定。

\paragraph{问题一：速度--能耗非线性优化 KKT 验证结果}

问题一的 KKT 验证结果来自“16 调速KKT验证”工作表。由于问题一共有 $1024$ 行主结果，且每一行分别检查二臂和三臂两个模式，因此 KKT 验证结果共有
\[
1024\times 2=2048
\]
行。其整体统计结果如表 \ref{tab:problem1_kkt_overall} 所示。

\begin{table}[H]
	\centering
	\caption{问题一速度--能耗非线性优化 KKT 整体验证结果}
	\label{tab:problem1_kkt_overall}
	\footnotesize
	\begin{tabular}{lcccc}
		\hline
		数据范围 & KKT检查行数 & 通过行数 & 未通过行数 & 通过率 \\
		\hline
		全部参数组 & 2048 & 2045 & 3 & 99.85\% \\
		基准参数组 & 512 & 512 & 0 & 100.00\% \\
		\hline
	\end{tabular}
\end{table}

为了说明 KKT 验证并不是只给出最终通过率，表 \ref{tab:problem1_kkt_samples} 给出若干典型检查行。表中“活跃约束”用于说明最优点附近主要由哪类约束起作用。

\begin{table}[H]
	\centering
	\caption{问题一速度--能耗非线性优化 KKT 验证样例}
	\label{tab:problem1_kkt_samples}
	\scriptsize
	\setlength{\tabcolsep}{3pt}
	\resizebox{\textwidth}{!}{
		\begin{tabular}{lccccccl}
			\hline
			任务组合 & 截止比例 & 模式 & 最优速度 $(s_S,s_D)$ & 优化能耗 & 优化后完工时间 & 活跃约束 & KKT是否通过 \\
			\hline
			\texttt{1122} & 0.80 & 二臂 & $(1.2436,1.2532)$ & 14704.08 & 2291.73 & 截止时间约束 & 是 \\
			\texttt{1122} & 0.80 & 三臂 & $(0.9163,0.9687)$ & 14464.63 & 2285.60 & 无 & 是 \\
			\texttt{2111} & 0.75 & 二臂 & $(1.3325,1.3346)$ & 8669.77 & 1466.50 & 截止时间约束 & 是 \\
			\texttt{2111} & 0.75 & 三臂 & $(0.9163,0.9687)$ & 8428.68 & 1391.63 & 无 & 是 \\
			\hline
		\end{tabular}
	}
\end{table}

由表 \ref{tab:problem1_kkt_overall} 和表 \ref{tab:problem1_kkt_samples} 可知，问题一的数值解整体稳定。数结果不仅满足速度上下界和截止时间约束，而且能够通过 KKT 一阶最优性检查。。

\paragraph{问题二：鲁棒速度模式选择输出结果}

问题二不再对多个 seed 取平均，而是要求同一任务数量组合下所有 seed 共用同一组速度倍率，并最小化最坏情形能耗。该问题共考虑 $16$ 种任务数量组合和 $16$ 个截止时间比例，因此主结果共有
\[
16\times 16=256
\]
行。其整体输出结果如表 \ref{tab:problem2_output_overall} 所示。

\begin{table}[H]
	\centering
	\caption{问题二鲁棒速度模式选择整体输出统计}
	\label{tab:problem2_output_overall}
	\footnotesize
	\setlength{\tabcolsep}{4pt}
	\begin{tabular}{lccccc}
		\hline
		模式 & 鲁棒可行组数 & 推荐次数 & 平均最坏能耗 & 单臂速度倍率 & 双臂速度倍率 \\
		\hline
		二臂模式 & 256 & 187 & 12404.60 & 1.0768 & 1.1196 \\
		三臂模式 & 256 & 69 & 12500.64 & 0.9247 & 0.9766 \\
		\hline
	\end{tabular}
\end{table}

为了展示鲁棒问题中推荐模式如何随截止时间变化，表 \ref{tab:problem2_output_samples} 选取了 \texttt{1111}、\texttt{1122} 和 \texttt{2222} 三组任务组合的代表性结果。表中 $Z^{2B}$ 和 $Z^{3B}$ 表示二臂、三臂模式下的最坏情形能耗上界。

\begin{table}[H]
	\centering
	\caption{问题二鲁棒速度模式选择代表性输出样例}
	\label{tab:problem2_output_samples}
	\tiny
	\setlength{\tabcolsep}{2.5pt}
	\resizebox{\textwidth}{!}{
		\begin{tabular}{cccccccccc}
			\hline
			任务组合 & 截止比例 & 截止时间 
			& 二臂速度 $(s_S,s_D)$ & $Z^{2B}$ & 二臂最大完工时间 
			& 三臂速度 $(s_S,s_D)$ & $Z^{3B}$ & 三臂最大完工时间 
			& 推荐模式 \\
			\hline
			\texttt{1111} & 0.75 & 1398.00
			& $(1.4031,1.4254)$ & 9658.23 & 1398.00
			& $(0.9163,0.9687)$ & 9055.26 & 1333.67
			& 三臂 \\
			\texttt{1111} & 0.80 & 1491.20
			& $(1.3115,1.3403)$ & 9324.95 & 1491.20
			& $(0.9163,0.9687)$ & 9055.26 & 1333.67
			& 三臂 \\
			\texttt{1111} & 0.85 & 1584.40
			& $(1.2313,1.2648)$ & 9076.26 & 1584.40
			& $(0.9163,0.9687)$ & 9055.26 & 1333.67
			& 三臂 \\
			\texttt{1122} & 0.85 & 2434.97
			& $(1.2173,1.2514)$ & 14983.77 & 2434.97
			& $(0.9163,0.9687)$ & 14785.66 & 2300.69
			& 三臂 \\
			\texttt{1122} & 0.90 & 2578.20
			& $(1.1454,1.1842)$ & 14694.07 & 2578.20
			& $(0.9163,0.9687)$ & 14785.66 & 2300.69
			& 二臂 \\
			\texttt{1122} & 0.95 & 2721.43
			& $(1.0814,1.1239)$ & 14491.85 & 2721.43
			& $(0.9163,0.9687)$ & 14785.66 & 2300.69
			& 二臂 \\
			\texttt{2222} & 0.70 & 2187.50
			& $(1.4835,1.5017)$ & 17747.13 & 2187.50
			& $(1.0603,1.1038)$ & 16051.60 & 2187.50
			& 三臂 \\
			\texttt{2222} & 0.75 & 2343.75
			& $(1.3809,1.4055)$ & 16996.60 & 2343.75
			& $(0.9859,1.0342)$ & 15899.76 & 2343.75
			& 三臂 \\
			\texttt{2222} & 0.80 & 2500.00
			& $(1.2912,1.3211)$ & 16435.10 & 2500.00
			& $(0.9211,0.9731)$ & 15850.69 & 2500.00
			& 三臂 \\
			\hline
		\end{tabular}
	}
\end{table}

由表 \ref{tab:problem2_output_samples} 可以看出，鲁棒问题的推荐结果与截止时间比例密切相关。在 \texttt{1122} 组合中，当截止比例为 $0.85$ 时，三臂最坏能耗低于二臂，因此推荐三臂；当截止比例放宽到 $0.90$ 和 $0.95$ 时，二臂最坏能耗反而更低，因此推荐二臂。这说明鲁棒问题并不是简单地因为三臂资源更多就总是推荐三臂，而是在紧工期下三臂更有优势，在工期逐渐宽松后二臂的低能耗优势会重新体现出来。

\paragraph{问题二：鲁棒速度模式选择 KKT 验证结果}

由于问题二共有 $256$ 行主结果，且每一行分别检查二臂和三臂两个模式，因此 KKT 验证结果共有
\[
256\times 2=512
\]
行。其整体统计结果如表 \ref{tab:problem2_kkt_overall} 所示。

\begin{table}[H]
	\centering
	\caption{问题二鲁棒速度模式选择 KKT 整体验证结果}
	\label{tab:problem2_kkt_overall}
	\footnotesize
	\begin{tabular}{lcccc}
		\hline
		检查对象 & KKT检查行数 & 通过行数 & 未通过行数 & 通过率 \\
		\hline
		二臂模式 & 256 & 107 & 149 & 41.80\% \\
		三臂模式 & 256 & 224 & 32 & 87.50\% \\
		合计 & 512 & 331 & 181 & 64.65\% \\
		\hline
	\end{tabular}
\end{table}

为了说明鲁棒 KKT 验证并非全部失败，表 \ref{tab:problem2_kkt_samples} 选取了若干代表性检查行，既包含 KKT 未通过的紧工期结果，也包含 KKT 通过的结果。表中可以看出，所有样例均满足鲁棒可行性，但是否满足严格 KKT 检查与活跃约束结构和平稳性残差有关。

\begin{table}[H]
	\centering
	\caption{问题二鲁棒速度模式选择 KKT 验证样例}
	\label{tab:problem2_kkt_samples}
	\tiny
	\setlength{\tabcolsep}{2.2pt}
	\resizebox{\textwidth}{!}{
		\begin{tabular}{ccccccccc}
			\hline
			任务组合 & 截止比例 & 模式 
			& 最优速度 $(s_S,s_D)$ 
			& 最坏能耗 $z$ 
			& 调速后最大完工时间 
			& 活跃约束 
			& 平稳性残差范数 
			& KKT是否通过 \\
			\hline
			\texttt{1111} & 0.70 & 二臂 
			& $(1.5070,1.5234)$ 
			& 10101.39 
			& 1304.80 
			& 时间种子1约束+能耗种子2约束 
			& 0.999967 
			& 否 \\
			
			\texttt{1111} & 0.70 & 三臂 
			& $(0.9378,0.9887)$ 
			& 9058.01 
			& 1304.80 
			& 时间种子2约束+能耗种子2约束 
			& 0.995768 
			& 否 \\
			
			\texttt{1111} & 0.75 & 三臂 
			& $(0.9163,0.9687)$ 
			& 9055.26 
			& 1333.67 
			& 能耗种子2约束 
			& 0.000000 
			& 是 \\
			
			\texttt{1111} & 1.15 & 二臂 
			& $(0.9163,0.9687)$ 
			& 8597.46 
			& 2099.25 
			& 能耗种子2约束 
			& 0.000000 
			& 是 \\
			
			\texttt{1122} & 0.80 & 三臂 
			& $(0.9203,0.9723)$ 
			& 14785.81 
			& 2291.70 
			& 时间种子2约束+能耗种子2约束 
			& 0.810700 
			& 否 \\
			
			\texttt{1122} & 0.85 & 三臂 
			& $(0.9163,0.9687)$ 
			& 14785.66 
			& 2300.69 
			& 能耗种子2约束 
			& 0.000000 
			& 是 \\
			
			\texttt{2211} & 0.85 & 三臂 
			& $(0.9163,0.9687)$ 
			& 10076.15 
			& 1739.35 
			& 能耗种子0约束 
			& 0.000000 
			& 是 \\
			
			\texttt{2222} & 0.80 & 三臂 
			& $(0.9211,0.9731)$ 
			& 15850.69 
			& 2500.00 
			& 时间种子0约束+能耗种子0约束 
			& 0.947640 
			& 否 \\
			
			\texttt{2222} & 0.85 & 三臂 
			& $(0.9163,0.9687)$ 
			& 15850.45 
			& 2512.31 
			& 能耗种子0约束 
			& 0.000000 
			& 是 \\
			\hline
		\end{tabular}
	}
\end{table}


由表 \ref{tab:problem2_kkt_overall} 和表 \ref{tab:problem2_kkt_samples} 可以看出，鲁棒问题的 KKT 通过率明显低于问题一。这并不说明鲁棒结果没有意义，而是说明逐次线性规划法得到的结果更适合解释为鲁棒可行解或数值近似解。由于鲁棒模型同时包含多个 seed 的时间约束和能耗上界约束，活跃约束集合可能在不同 seed 之间切换，乘子恢复和平稳性检查更敏感，因此 KKT 通过率低于问题一是合理的。


\subsection{解的评估与结果解释}

对于主问题，由于 CP-SAT 求解器是成熟的标准求解器，因此本文将 CP-SAT 结果作为对照组，比较自编算法与其在最大完工时间、系统总能耗和求解时间上的差异。若自编方法得到的目标值与 CP-SAT 保持一致，则说明该方法在对应算例上达到了标准求解器结果；若存在差异，则进一步分析差异是否来自启发式近似或目标优先级差异。

对于非线性延伸问题，通过 KKT 条件验证数值解的可靠性。问题一为速度--能耗非线性优化问题，目标函数和约束形式较规整，KKT 验证通过率较高；问题二为鲁棒速度模式选择问题，涉及多个 seed 的最坏情形约束，采用逐次线性规划近似求解，因此 KKT 通过率较低，需要区分鲁棒可行解和严格 KKT 最优解。

\subsubsection{主问题解的评估}

主问题的解主要从四种求解方法之间的对照关系进行评价。本文将 CP-SAT 求解器作为标准对照组，因为 CP-SAT 能够直接处理 0--1 变量、可选区间变量、资源互斥约束和多阶段目标，在本文主问题中可作为判断其他方法结果是否合理的基准。

从上述数据可以看出，CP-SAT 求解器、融合剪枝的隐枚举法和热启动后的隐枚举法在二臂平均完工时间、二臂平均能耗、三臂平均完工时间和三臂平均能耗上完全一致。具体来说，三种方法均得到二臂平均完工时间 $2554.17$、二臂平均能耗 $11479.25$、三臂平均完工时间 $1758.69$ 和三臂平均能耗 $11976.56$。这说明融合剪枝的隐枚举法和热启动后的隐枚举法在这些算例上达到了与 CP-SAT 相同的目标值，其求解结果是合理的。

该一致性也说明，本文自编精确搜索方法中的剪枝并没有错误删除最优分支。若剪枝规则过强，可能出现自编方法的结果优于或劣于 CP-SAT 的异常情况；而现在三种方法的完工时间和能耗完全一致，说明完工时间下界、能耗下界、负载下界和中心安全区互斥剪枝能够减少搜索量，同时保留最优解。因此，融合剪枝的隐枚举法不仅在逻辑上符合分支定界思想，在数值结果上也得到了标准求解器的验证。

启发式算法的结果与 CP-SAT 并不完全一致，但这种差异是合理的。启发式算法的二臂平均完工时间为 $2555.23$，略高于 CP-SAT 的 $2554.17$；三臂平均完工时间为 $1781.65$，高于 CP-SAT 的 $1758.69$。这说明启发式算法没有在所有算例上达到全局最优，尤其在三臂场景中偏差更明显。

这种偏差并不表示启发式结果无效，而是由启发式方法的性质决定的。启发式算法主要依靠规则构造和局部改进，目标是在较短时间内生成质量较好的可行解，而不是保证遍历完整搜索空间。三臂系统中双臂协同任务的可选组合更多，任务顺序和机械臂组合之间的耦合更强，因此启发式规则更容易陷入局部较优方案。表 \ref{tab:heuristic_gap_summary} 中二臂模式有 $44/48$ 组达到最优，而三臂模式只有 $23/48$ 组达到最优，也说明三臂场景的搜索难度更高。

\subsubsection{非线性延伸问题解的评估}

非线性延伸问题的解主要通过 KKT 条件进行评价。KKT 条件包括原问题可行性、乘子非负性、互补松弛和平稳性条件。对于本文的凸优化问题，在满足一定正则性条件的情况下，KKT 条件不仅是局部最优的一阶必要条件，也可以作为判断数值解是否具有最优性的重要依据。因此，本文在非线性部分不只检查解是否满足约束，还进一步检查其是否满足 KKT 条件。

对于问题一，即速度--能耗非线性优化问题，KKT 验证结果整体较好。由表 \ref{tab:problem1_kkt_overall} 可知，全部参数组下共有 $2048$ 行 KKT 检查结果，其中 $2045$ 行通过，未通过仅 $3$ 行，通过率达到 $99.85\%$；基准参数组下 $512$ 行全部通过，通过率为 $100.00\%$。

该结果说明，内点制约函数法结合 $0.618$ 一维搜索得到的速度解，在绝大多数情况下不仅满足速度上下界和截止时间约束，也满足互补松弛和平稳性检查。由于问题一只有两个主要连续变量 $s_S$ 和 $s_D$，且时间函数由 $1/s$ 项构成，能耗函数由 $1/s$ 和 $s^2$ 项构成，在正速度区间内具有较好的凸性结构，因此 KKT 检查通过率高是合理的。

对于问题一中少数未通过 KKT 的结果，不应理解为模型错误，而更可能来自数值迭代精度、边界附近残差或容差设置。因为其未通过数量仅占全部结果的极小比例，并且总体推荐趋势和能耗比较不依赖这几行异常结果，所以问题一的数值解可以认为整体稳定，能够支持后续关于二臂/三臂速度--能耗权衡的分析。

对于问题二，即鲁棒速度模式选择问题，KKT 验证结果明显复杂。由表 \ref{tab:problem2_kkt_overall} 可知，鲁棒问题共有 $512$ 行 KKT 检查结果，其中通过 $331$ 行，未通过 $181$ 行，总通过率为 $64.65\%$。其中，二臂模式通过率为 $41.80\%$，三臂模式通过率为 $87.50\%$。

问题二 KKT 通过率较低是可以解释的。
\begin{enumerate}
	\item 问题二不是对平均结果进行优化，而是要求同一组速度同时满足多个 seed 的截止时间约束和能耗上界约束。也就是说，模型中存在多组时间约束和多组能耗约束，最优点可能由某一个或多个最不利 seed 决定。当最坏能耗 seed 或活跃时间约束在迭代过程中发生切换时，活跃约束集合会变得敏感，从而导致 KKT 平稳性残差较大。

	\item 采用逐次线性规划法求解。该方法每一步求解的是当前点附近的线性化子问题，而不是一次性直接求解原始非线性模型。虽然程序会将候选解代回原非线性约束中检查可行性，但最终结果仍可能只是逐次线性化意义下的近似收敛点。因此，某些结果能够满足鲁棒可行性，却未必满足原始非线性问题的严格 KKT 条件。

	\item 二臂模式 KKT 通过率低于三臂模式也是合理的。二臂系统资源较少，为了让所有 seed 都满足同一截止时间，往往需要更高速度，并且更容易使时间约束处于活跃状态。当时间约束和能耗上界约束同时接近边界时，乘子恢复和平稳性检查难度会增加，因此二臂模式更容易出现 KKT 未通过。三臂系统由于并行能力更强，时间压力相对较小，活跃约束结构更稳定，所以 KKT 通过率更高。
\end{enumerate}

因此，对问题二结果的解释应更加谨慎。KKT 未通过并不表示该行结果不可行，也不表示鲁棒模型没有意义；它只说明该结果不能被严格称为满足一阶最优性的局部最优解。对于 KKT 未通过的结果，应表述为“鲁棒可行解”或“数值近似解”；对于 KKT 通过的结果，则可以进一步认为其满足一阶最优性检查。

从工程分析角度看，问题二仍然有意义。因为鲁棒模型的首要目标是找到一组统一速度，使同一任务数量组合下所有 seed 都满足截止时间约束，并尽量控制最坏情形能耗。即使部分结果未通过严格 KKT 检查，只要它们通过原非线性约束的可行性检查，仍然可以用于比较二臂和三臂在 seed 扰动下的稳健性差异。报告中应避免把这些结果表述为严格最优解，而应表述为鲁棒可行的数值近似结果。

综上，问题一的 KKT 通过率接近 $100\%$，说明内点法与 $0.618$ 搜索得到的速度--能耗优化结果整体可靠；问题二由于采用逐次线性规划法，并且包含多个 seed 的最坏情形约束，KKT 通过率较低，因此其结果更适合解释为鲁棒可行解和数值近似解，而非严格最优解。


\subsection{方法对比}

\subsubsection{主问题四种方法对比}


四种方法在 16 种任务组合下的最优性与平均计算时间如表 \ref{tab:main_method_opt_time} 所示。其中，计算时间为同一任务组合下不同 seed 和不同机械臂模式的平均求解时间。

\begin{table}[H]
	\centering
	\caption{主问题四种方法的最优性与计算时间对比}
	\label{tab:main_method_opt_time}
	\tiny
	\setlength{\tabcolsep}{2.2pt}
	\resizebox{\textwidth}{!}{
		\begin{tabular}{ccccccccc}
			\hline
			任务组合 
			& 方法1是否最优 & 方法1时间/s
			& 方法2是否最优 & 方法2时间/s
			& 方法3是否最优 & 方法3时间/s
			& 方法4是否最优 & 方法4时间/s \\
			\hline
			\texttt{1111} & 是 & 1.1087 & 是 & 0.0160 & 是 & 0.1873 & 是 & 0.0374 \\
			\texttt{1112} & 是 & 0.4266 & 是 & 0.1280 & 是 & 0.2025 & 是 & 0.1635 \\
			\texttt{1121} & 是 & 0.3741 & 是 & 0.1290 & 否 & 0.2024 & 是 & 0.1564 \\
			\texttt{1122} & 是 & 2.8354 & 是 & 1.7181 & 是 & 1.0874 & 是 & 1.6760 \\
			\texttt{1211} & 是 & 0.3500 & 是 & 0.1129 & 否 & 0.2134 & 是 & 0.1565 \\
			\texttt{1212} & 是 & 0.9333 & 是 & 0.5551 & 否 & 1.0964 & 是 & 0.6132 \\
			\texttt{1221} & 是 & 1.6254 & 是 & 0.7532 & 否 & 1.1362 & 是 & 0.7846 \\
			\texttt{1222} & 是 & 17.3515 & 是 & 18.3875 & 否 & 1.0779 & 是 & 9.8419 \\
			\texttt{2111} & 是 & 0.3118 & 是 & 0.1284 & 否 & 0.2060 & 是 & 0.1477 \\
			\texttt{2112} & 是 & 0.9109 & 是 & 0.7779 & 否 & 1.0458 & 是 & 0.7320 \\
			\texttt{2121} & 是 & 1.0879 & 是 & 0.8380 & 否 & 1.1137 & 是 & 0.8028 \\
			\texttt{2122} & 是 & 19.1113 & 是 & 27.9397 & 否 & 1.2413 & 是 & 23.5496 \\
			\texttt{2211} & 是 & 0.6799 & 是 & 1.2688 & 否 & 1.1127 & 是 & 0.5756 \\
			\texttt{2212} & 是 & 3.2290 & 是 & 3.3501 & 否 & 1.2333 & 是 & 4.1532 \\
			\texttt{2221} & 是 & 3.6314 & 是 & 4.2431 & 否 & 1.0568 & 是 & 4.7159 \\
			\texttt{2222} & 是 & 62.2421 & 是 & 94.0316 & 否 & 1.8600 & 是 & 48.9079 \\
			\hline

		\end{tabular}
	}
\end{table}

由表 \ref{tab:main_method_opt_time} 可以看出，方法 1、方法 2 和方法 4 在 16 种任务组合下均能够得到与 CP-SAT 一致的结果，因此这三种方法都可以认为能够找到主问题的最优解。方法 3 的计算时间最短，平均仅为 $0.8796$ s，但其结果并不总是与 CP-SAT 一致，因此启发式算法更适合作为快速可行解生成方法，而不适合作为严格最优解算法单独使用。

在能够找到最优解的三种方法中，方法 2 与方法 1 的速度关系随任务规模变化而变化。当任务数量较少、组合结构较简单时，方法 2 往往比方法 1 更快。例如在 \texttt{1111}、\texttt{1112}、\texttt{1121}、\texttt{1211} 等组合中，方法 2 的计算时间均低于方法 1。这说明在小规模算例中，自编隐枚举法结合剪枝可以较快排除无效分支，搜索效率较高。

但是，当任务数量增加，尤其是 Type3、Type4 等双臂协同任务数量增加时，方法 2 的搜索空间迅速扩大，计算时间开始超过方法 1。例如在 \texttt{1222} 组合中，方法 2 用时 $18.3875$ s，高于方法 1 的 $17.3515$ s；在 \texttt{2122} 组合中，方法 2 用时 $27.9397$ s，高于方法 1 的 $19.1113$ s；在 \texttt{2222} 组合中，方法 2 用时达到 $94.0316$ s，也高于方法 1 的 $62.2421$ s。这说明单纯依靠剪枝的隐枚举法虽然能够保证最优性，但在较大规模组合下仍会受到搜索空间膨胀的影响。

方法 4 在保持最优性的同时，整体计算时间最优。其平均计算时间为 $6.0634$ s，低于方法 1 的 $7.2631$ s 和方法 2 的 $9.6486$ s。尤其在较大规模组合下，方法 4 的优势更加明显。例如 \texttt{1222} 组合中，方法 4 用时 $9.8419$ s，明显低于方法 1 和方法 2；\texttt{2222} 组合中，方法 4 用时 $48.9079$ s，也低于方法 1 的 $62.2421$ s 和方法 2 的 $94.0316$ s。

这说明热启动策略是有效的。方法 4 先利用启发式算法快速得到一个较好的初始可行解，再以该解作为隐枚举搜索的初始上界，从而使大量不可能优于当前解的分支被提前剪掉。与方法 2 相比，方法 4 并没有改变隐枚举法的最优性判断逻辑，因此仍然能够得到最优解；但由于初始上界更紧，搜索效率明显提高。

综上，在四种方法中，方法 3 的速度最快但不能保证最优性；方法 1 可以作为可靠对照；方法 2 能够找到最优解，但在任务数量较多时速度下降明显；方法 4 在保持最优性的同时具有更好的总体计算效率。因此，本文将方法 4，即热启动后的隐枚举法，作为主问题中综合表现较好的算法。



\subsubsection{延伸问题方法对比}

表 \ref{tab:nonlinear_method_selection} 对课内常见方法与本文方法的关系进行对比。

\begin{table}[H]
	\centering
	\caption{延伸问题非线性规划方法选择对比}
	\label{tab:nonlinear_method_selection}
	\scriptsize
	\setlength{\tabcolsep}{3pt}
	\resizebox{\textwidth}{!}{
		\begin{tabular}{p{3.0cm}p{5.2cm}p{6.2cm}p{3.0cm}}
			\hline
			方法类别 & 课内方法特点 & 与本文问题的关系 & 本文处理方式 \\
			\hline
			解析法与 KKT 条件 
			& 通过一阶条件、二阶条件或 KKT 条件描述最优点应满足的数学性质 
			& 本文两个问题均含不等式约束，KKT 条件适合用于检验数值解是否满足可行性、互补松弛和平稳性 
			& 用于解的验证，而不是作为主要求解算法 \\
			
			最速下降法 
			& 根据负梯度方向迭代，计算简单，适合说明下降迭代思想 
			& 问题一、问题二都存在速度边界和截止时间约束，单纯沿负梯度方向迭代还需要额外处理可行性 
			& 吸收其“下降方向”思想，但不单独采用 \\
			
			牛顿法、拟牛顿法、共轭梯度法 
			& 适合无约束或经过转化后的光滑优化问题，强调利用梯度或近似 Hessian 信息提高收敛效率 
			& 本文变量维数较低，且约束条件比变量维数更重要，重点不是高维无约束搜索，而是保持约束可行并处理截止时间限制 
			& 不作为主要算法，主要保留梯度与一阶条件思想 \\
			
			0.618 法 
			& 用于一维搜索，在单峰区间内不断缩小最优步长所在范围 
			& 问题一在确定下降方向后，需要沿该方向选择合适步长；由于变量少、目标函数光滑，一维搜索结构清楚 
			& 与内点制约函数法结合，用于问题一的步长搜索 \\
			
			制约函数法 
			& 将有约束问题转化为带制约项的无约束或近似无约束问题，使迭代点保持在可行域内部 
			& 问题一具有明确的不等式约束，包括速度上下界和截止时间约束，适合用制约函数统一处理 
			& 作为问题一的主要求解方法 \\
			
			逐次逼近法 
			& 将复杂非线性问题在当前点附近近似为较易求解的子问题，再逐步迭代逼近原问题解 
			& 问题二包含多个 seed 的时间约束和能耗上界约束，引入辅助变量后可在当前点附近线性化为小规模线性规划子问题 
			& 采用逐次线性规划法求解问题二 \\
			\hline
		\end{tabular}
	}
\end{table}

对于问题一，本文选择内点制约函数法结合 $0.618$ 法。问题一的变量只有单臂速度倍率 $s_S$ 和双臂协同速度倍率 $s_D$，约束主要是速度上下界和截止时间约束，均属于不等式约束：
\[
s_{\min}\leq s_S\leq s_{\max},\qquad
s_{\min}\leq s_D\leq s_{\max},
\]
\[
C_{\max}(s_S,s_D)\leq D.
\]
因此，使用制约函数法可以将约束统一写入目标函数，使迭代过程保持在可行域内部。随后，在每一步搜索中，二维问题可以沿下降方向转化为关于步长 $\alpha$ 的一维搜索问题：
\[
\varphi(\alpha)=\Phi_\mu(\boldsymbol{s}^{(k)}+\alpha\boldsymbol{d}^{(k)}).
\]
此时采用 $0.618$ 法只需要比较函数值，不需要计算二阶导数，适合本文变量少、函数形式明确的速度--能耗优化问题。

对于问题二，本文选择逐次线性规划法。问题二需要所有 seed 共用同一组速度倍率，并最小化最坏情形能耗。引入辅助变量 $z$ 后，模型可写为
\[
\min_{s_S,s_D,z}\quad z,
\]
\[
E^{(r)}(s_S,s_D)\leq z,\qquad
C_{\max}^{(r)}(s_S,s_D)\leq D,\qquad \forall r\in\mathcal{R}.
\]
该问题变量数量仍然较少，但约束数量随 seed 增加而增加。逐次线性规划法可以在当前点附近对非线性约束进行一阶线性化，将原问题逐步转化为小规模线性规划子问题，便于自编程序实现，也适合处理多 seed 鲁棒约束。

综上，本文的问题一制约函数法处理约束、$0.618$ 法确定步长，问题二采用逐次线性规划法处理多 seed 鲁棒约束。KKT 条件则用于求解后的最优性验证。这样的安排既对应了课内非线性规划中的制约函数法、步长搜索、逐次逼近法和 KKT 条件，也与本文两个非线性问题的结构相匹配。

